\documentclass[11pt,a4paper]{article}
\usepackage{fontspec}
\setmainfont{DejaVu Sans}[Scale=0.90]
\usepackage{amsmath,amssymb}
\usepackage[margin=2.4cm]{geometry}
\usepackage{booktabs,array,longtable,tabularx}
\usepackage{enumitem}
\usepackage{titlesec}
\usepackage{parskip}
\titlelabel{\thetitle.\quad}
\titleformat{\section}{\normalfont\large\bfseries}{\thesection.}{0.6em}{}
\titleformat{\subsection}{\normalfont\normalsize\bfseries}{\thesubsection}{0.6em}{}
\titlespacing*{\section}{0pt}{16pt}{6pt}
\titlespacing*{\subsection}{0pt}{12pt}{4pt}
\setlist[itemize]{leftmargin=1.4em,itemsep=1pt,topsep=3pt}
\newcolumntype{L}[1]{>{\raggedright\arraybackslash}p{#1}}
\begin{document}
\begin{center}
{\Large\bfseries Apport et limites --- positionnement}\\[4pt]
{\small Krapez, \emph{International Journal of Thermal Sciences} \textbf{136} (2018) 182--199}\\[2pt]
{\small Version corrigée --- 8 septembre 2026}
\end{center}

space{4pt}


\section{Ce que l'article établit}

\subsection{Résultats acquis}

\begin{center}
\begin{tabularx}{\textwidth}{@{}lX@{}}
\toprule
Transformation de Liouville & Établie pour l'équation de la chaleur en milieu à gradient, sous les deux
formulations duales \\
\addlinespace
Identification du potentiel & $V = s''/s$, courbure relative de la métapropriété, elle-même racine de
l'effusivité ou son inverse \\
\addlinespace
Réduction du nombre d'inconnues & De deux fonctions matériau à une seule, l'effusivité exprimée en
fonction de la coordonnée transformée \\
\addlinespace
Trois familles solvables & Issues d'un potentiel constant, avec solutions en fonctions élémentaires
exclusivement \\
\addlinespace
Formulations quadripolaires & Établies pour les trois familles, sous les deux formes duales \\
\addlinespace
Indétermination du retour & Un profil transformé correspond à une infinité de profils dans l'espace
physique \\
\bottomrule
\end{tabularx}
\end{center}

\subsection{Portée méthodologique}

La disponibilité de quadripôles gradués en fonctions élémentaires permet de représenter un profil réel
avec un nombre de couches inférieur à celui qu'exige le modèle classique par morceaux constants. Le
gain porte sur le coût de calcul et sur le nombre de paramètres d'ajustement.

Le paramètre $\xi_{c}$ des familles hyperbolique et trigonométrique est non linéaire, ce qui impose un
solveur non linéaire lors de l'ajustement. Les paramètres de la famille linéaire, eux, interviennent
linéairement.

\section{Ce que l'article n'aborde pas}

\subsection{Cadre constitutif}

L'article se place intégralement sous la loi de Fourier. Le cas d'une loi à temps de relaxation fini
n'est pas examiné, et le paramètre spectral y demeure purement imaginaire en régime harmonique.

\subsection{Lecture spectrale}

L'auteur écarte explicitement l'interprétation spectrale de l'opérateur obtenu. La forme de Schrödinger
est employée comme instrument de construction de solutions exactes, non comme objet d'analyse. Aucune
recherche de valeurs propres, d'états liés ou d'états de diffusion n'est entreprise, et la position du
paramètre $E = -p$ dans le plan complexe n'est pas discutée.

\subsection{Conditionnement}

L'article constate la restitution imparfaite des variations profondes d'effusivité et la relie à la
difficulté de l'inversion thermique en profondeur. Ce constat demeure phénoménologique : il n'est
rattaché ni à la structure de l'opérateur, ni à un critère quantitatif portant sur la bande de mesure.

\section{Injectivité et conditionnement}

\subsection{Deux énoncés distincts}

L'article contient deux résultats qui peuvent paraître contradictoires. D'une part, deux profils duaux
distincts ne peuvent jamais produire la même réponse thermique à toute fréquence : l'application
directe est injective. D'autre part, des différences importantes dans les propriétés profondes
demeurent pratiquement invisibles dans le signal de surface.

La contradiction n'est qu'apparente. L'injectivité est une propriété de l'application mathématique
considérée sur l'ensemble du domaine fréquentiel et en l'absence de bruit. Le conditionnement mesure la
sensibilité de la réponse aux variations du profil, sur une bande de mesure finie et à niveau de bruit
donné.

\subsection{Formulation}

Un problème inverse peut être injectif tout en étant pratiquement indéterminé. L'unicité garantit qu'il
existe au plus une solution ; elle ne garantit ni sa stabilité, ni la possibilité de l'atteindre à
partir de données bruitées sur une bande finie.

Cette distinction sépare deux énoncés fréquemment confondus dans la littérature de métrologie
thermique, où l'expression « problème mal posé » recouvre indifféremment la non-unicité et le mauvais
conditionnement.

\subsection{Deux sources distinctes de non-détermination}

La non-unicité structurelle établie à la section 2.4 de l'article, associée à l'indétermination de la
transformation inverse, est indépendante du bruit et de la bande de mesure. Elle subsiste avec des
données parfaites.

Le mauvais conditionnement, lui, dépend de la bande de fréquences accessible et du niveau de bruit. Il
s'ajoute à la première sans s'y réduire.

\section{Points ouverts}

\begin{center}
\begin{tabularx}{\textwidth}{@{}lXX@{}}
\toprule
\textbf{Question} & \textbf{État dans l'article} & \textbf{Reste à établir} \\
\midrule
Position du paramètre spectral & Sur l'axe imaginaire, non discutée &
Comportement sous une loi constitutive à temps de relaxation fini \\
\addlinespace
Origine du mauvais conditionnement & Constatée, non expliquée &
Rattachement à la structure spectrale de l'opérateur \\
\addlinespace
Critère quantitatif d'identifiabilité & Absent &
Condition portant sur la bande de mesure et le niveau de bruit \\
\addlinespace
Transposition à l'inversion & Annoncée comme travaux en cours (2018) &
Vérification des publications ultérieures de l'auteur \\
\addlinespace
Perte d'unicité en régime ondulatoire & Rapportée pour le cas électromagnétique &
Portée pour un régime intermédiaire entre diffusion et propagation \\
\bottomrule
\end{tabularx}
\end{center}

\section{Antériorités à considérer}

\begin{center}
\begin{tabularx}{\textwidth}{@{}lX@{}}
\toprule
Krapez (2016) & \emph{Int. J. Heat Mass Transf.} \textbf{99}, 485--503. Méthode PROFIDT et
transformations de Darboux conjointes ; travail fondateur du présent article \\
\addlinespace
Krapez, postérieur à 2018 & À rechercher : l'auteur annonce en conclusion la mise en {\oe}uvre de ces
outils pour les problèmes d'inversion \\
\addlinespace
Chen, Gamba, Li \& Wang (2025) & arXiv:2502.19533, accepté au \emph{SIAM J. Appl. Math.}
Reconstruction numérique du temps de relaxation par optimisation sous contrainte \\
\addlinespace
Camacho de la Rosa \emph{et al.} (2025) & \emph{J. Appl. Phys.} \textbf{137}, 155103. Caractérisation
des temps de relaxation par thermoréflectance en domaine fréquentiel \\
\addlinespace
Hoque \emph{et al.} (2024) & \emph{Appl. Phys. Lett.} \textbf{125}, 262201. Transition balistique vers
diffusif mesurée sur films d'AlN de 1,6 à 2440 nm \\
\addlinespace
Maillet \emph{et al.} (2000) & \emph{Thermal Quadrupoles}, Wiley. Ouvrage de référence du formalisme
matriciel employé \\
\bottomrule
\end{tabularx}
\end{center}

\section{Positionnement du présent travail}

Le cadre de l'article est celui de la loi de Fourier, sous lequel le paramètre spectral s'écrit
$p = i\omega$ et l'énergie associée à la forme de Schrödinger reste sur l'axe imaginaire. Dans ce
régime, la structure spectrale --- états liés, seuils, résonances --- est sans objet, ce qui justifie
que l'auteur écarte cette lecture.

Sous une loi constitutive à temps de relaxation fini, le paramètre devient $p + \tau p^{2}$, soit
$i\omega - \tau\omega^{2}$ en régime harmonique. L'énergie
\[ E = \tau\omega^{2} - i\omega \]
acquiert une partie réelle, et son argument
\[ \arg(E) = -\arctan\!\left(\frac{1}{\omega\tau}\right) \]
ne dépend que du produit $\omega\tau$. Il vaut $-90^{\circ}$ dans la limite diffusive, $-45^{\circ}$
exactement pour $\omega\tau = 1$, et tend vers $0^{\circ}$ dans la limite opposée, où l'énergie approche
l'axe réel positif.

Le prolongement envisagé porte donc sur la conséquence de ce déplacement pour le conditionnement du
problème inverse. Il s'appuie sur la transformation établie par Krapez, qu'il ne reproduit pas, et se
distingue des travaux de Chen \emph{et al.}, qui traitent la reconstruction numérique du temps de
relaxation sans analyse de la sensibilité en fonction de la bande de mesure.

Le résultat rapporté par l'auteur en section 4, selon lequel le passage à l'équation d'onde
s'accompagne d'une perte d'unicité dans le cas électromagnétique, constitue l'objection principale à
cette orientation et devra être traité explicitement.
\end{document}